\chapter{Variables and Expansion}\label{ch:mk-vars}

Make has variables for the same reason every language does, and one extra: they
are how you stop writing \cmd{gcc -Wall -Wextra -g} in four places and
discovering later that one of them says \cmd{-Wall} only.

\section{Recursive versus Simple Assignment}\label{sec:mk-assign}

There are two assignment operators and they are not interchangeable.

\begin{keys}[0.20]
	\texttt{=}   & \vocab{Recursive}. The right-hand side is stored unexpanded and re-expanded at every use \\
	\texttt{:=}  & \vocab{Simple}. The right-hand side is expanded once, now, and the result is stored \\
\end{keys}

\begin{makecode}
# recursive: expanded when used
CFLAGS  = -Wall $(EXTRA)
EXTRA   = -O2
# $(CFLAGS) is "-Wall -O2" -- EXTRA was defined later, and that is fine

# simple: expanded here and now
SOURCES := $(wildcard *.c)
# the directory is scanned once, not once per reference
\end{makecode}

\begin{note}[The trap]
	Recursive assignment re-expands \emph{every time}. Writing
	\cmd{SOURCES = \$(shell find ...)} with \texttt{=} therefore runs \cmd{find}
	again at every reference, possibly dozens of times in one build. Worse, a recursive
	variable that refers to itself is a fatal error:
\begin{makecode}
CFLAGS = $(CFLAGS) -g     # *** Recursive variable references itself.
CFLAGS := $(CFLAGS) -g    # fine: the right side is the OLD value
\end{makecode}
\end{note}

\begin{moral}
	Use \texttt{:=} by default. Reach for \texttt{=} only when you genuinely need
	late binding --- when the right-hand side mentions a variable that a later
	target-specific assignment will change.
\end{moral}

\section{Conditional and Appending Assignment}\label{sec:mk-condassign}

\begin{keys}[0.20]
	\texttt{?=}  & Assign only if the variable is not already set. Lets a caller override without editing the file \\
	\texttt{+=}  & Append, keeping the flavour of the original assignment          \\
	\texttt{!=}  & Assign the output of a shell command (GNU make 4.0 and later)    \\
\end{keys}

\begin{makecode}
CC      ?= gcc                  # `make CC=clang' wins; otherwise gcc
PREFIX  ?= /usr/local

CFLAGS  := -Wall -Wextra
CFLAGS  += -std=c17             # append
CFLAGS  += -g                   # append again

GIT_REV != git rev-parse --short HEAD
\end{makecode}

\begin{note}
	\texttt{?=} is how a Makefile stays configurable. A user building your project
	can write \cmd{make CC=clang CFLAGS=-O3} and never touch the file --- and
	packagers, who build thousands of projects, rely on exactly this.
\end{note}

\section{Automatic Variables}\label{sec:mk-auto}

Inside a recipe, make defines a handful of variables describing the rule it is
currently running. They are what make pattern rules possible at all, and they
are the least readable part of the language.

\begin{keys}[0.16]
	\texttt{\$@}  & The target                                                    \\
	\texttt{\$<}  & The \emph{first} prerequisite                                 \\
	\texttt{\$\textasciicircum{}} & \emph{All} prerequisites, duplicates removed  \\
	\texttt{\$+}  & All prerequisites, duplicates kept, original order            \\
	\texttt{\$?}  & Only the prerequisites newer than the target                  \\
	\texttt{\$*}  & The \vocab{stem}: what \texttt{\%} matched (\autoref{sec:mk-pattern}) \\
	\texttt{\$(@D)}, \texttt{\$(@F)} & The target's directory and filename parts  \\
\end{keys}

\begin{eg}
	The same rule, spelled out and then with automatic variables:
\begin{makecode}
app: main.o util.o
    gcc -o app main.o util.o

app: main.o util.o
    gcc -o $@ $^
\end{makecode}
	Read it as \emph{``make the target out of all the prerequisites''}. The second
	form cannot fall out of step with its own prerequisite list, which is the real
	argument for it.
\end{eg}

\begin{note}
	\texttt{\$<} and \texttt{\$\textasciicircum{}} differ in exactly the case that
	matters. Compiling one source into one object wants \texttt{\$<}, because the
	prerequisites include headers and \cmd{gcc -c} must not be handed
	\file{parser.h}. Linking wants \texttt{\$\textasciicircum{}}, because every
	object belongs on the command line.
\begin{makecode}
main.o: main.c parser.h
    gcc -c -o $@ $<            # just main.c
app: main.o parser.o
    gcc -o $@ $^               # both objects
\end{makecode}
\end{note}

\section{Functions and Substitution References}\label{sec:mk-funcs}

Make has a small functional language for manipulating lists of words. A call is
written as a dollar sign, a parenthesis, the function name, and its arguments;
the arguments are almost always filenames.

\begin{keys}[0.32]
	\texttt{\$(wildcard *.c)}            & Expand a glob against the filesystem      \\
	\texttt{\$(patsubst \%.c,\%.o,\$(SRC))} & Pattern substitution over a list       \\
	\texttt{\$(SRC:.c=.o)}               & The same thing, in shorthand              \\
	\texttt{\$(subst a,b,\$(TXT))}       & Plain textual substitution                \\
	\texttt{\$(filter \%.c,\$(FILES))}   & Keep the words that match                 \\
	\texttt{\$(filter-out main.o,\$(OBJ))} & Drop the words that match               \\
	\texttt{\$(sort \$(LIST))}           & Sort, and remove duplicates               \\
	\texttt{\$(words \$(LIST))}          & How many words                            \\
	\texttt{\$(dir \$(F))}, \texttt{\$(notdir \$(F))} & Directory and filename parts \\
	\texttt{\$(basename \$(F))}, \texttt{\$(suffix \$(F))} & Without and only the extension \\
	\texttt{\$(addprefix build/,\$(OBJ))} & Prepend to every word                    \\
	\texttt{\$(shell date +\%F)}         & Run a shell command, take its output      \\
	\texttt{\$(foreach v,\$(LIST),\dots)} & Loop over a list                         \\
\end{keys}

\begin{eg}[The idiom you will use every time]
\begin{makecode}
SRC := $(wildcard src/*.c)
OBJ := $(SRC:.c=.o)
# or, equivalently and more explicitly:
OBJ := $(patsubst %.c,%.o,$(SRC))
\end{makecode}
	Two lines, and the object list maintains itself. Adding a source file to
	\file{src/} needs no edit at all.
\end{eg}

\begin{note}
	\texttt{\$(wildcard)} sees the filesystem \emph{when the Makefile is read}. If
	a rule generates a \texttt{.c} file during the build, that file will not be in
	\texttt{\$(SRC)} --- make already made up its mind. Generated sources need
	their names written out, or a second invocation of make.
\end{note}

\begin{remark}
	Make's functions are whitespace-sensitive in a way that catches everyone once:
	\cmd{\$(filter \%.c, \$(FILES))} with a space after the comma passes
	\emph{``space plus the list''} as the second argument. Do not put spaces after
	the commas in a function call.
\end{remark}

\section{Conventional Variables}\label{sec:mk-conv}

Make's built-in rules (\autoref{sec:mk-implicit}) use a fixed set of variable
names, and so does every Makefile anyone else will read. Use these names and a
great deal works for free.

\begin{keys}[0.20]
	\texttt{CC}        & The C compiler. Defaults to \texttt{cc}                     \\
	\texttt{CXX}       & The C++ compiler                                            \\
	\texttt{CFLAGS}    & Flags for the C compiler                                    \\
	\texttt{CXXFLAGS}  & Flags for the C++ compiler                                  \\
	\texttt{CPPFLAGS}  & Preprocessor flags --- \cmd{-I}, \cmd{-D}. Used by both     \\
	\texttt{LDFLAGS}   & Linker flags --- \cmd{-L}, \cmd{-Wl,\dots}                  \\
	\texttt{LDLIBS}    & Libraries to link --- \cmd{-lm}, \cmd{-lpthread}            \\
	\texttt{ARFLAGS}   & Flags for \cmd{ar}                                          \\
	\texttt{PREFIX}    & Where \texttt{install} should put things                    \\
\end{keys}

\begin{note}
	\texttt{LDFLAGS} and \texttt{LDLIBS} are not the same, and the split is not
	pedantry. \texttt{LDFLAGS} says \emph{where} to look and goes early on the
	command line; \texttt{LDLIBS} says \emph{what} to link and must come after the
	object files, because the GNU linker resolves symbols in the order it is given
	them. Putting \cmd{-lm} in \texttt{LDFLAGS} produces
	\texttt{undefined reference to `sqrt'} on a program that plainly links the
	maths library.
\end{note}

\begin{makecode}
CC       := gcc
CPPFLAGS := -Iinclude -DNDEBUG
CFLAGS   := -Wall -Wextra -std=c17 -O2
LDFLAGS  := -Llib
LDLIBS   := -lm

app: $(OBJ)
    $(CC) $(LDFLAGS) -o $@ $^ $(LDLIBS)
\end{makecode}

\section{Overriding from Outside}\label{sec:mk-override}

Variables can be set in four places. Later entries win:

\begin{keys}[0.30]
	1. The Makefile           & \cmd{CFLAGS := -O2}                               \\
	2. The environment        & \cmd{CFLAGS=-O3 make} --- but only for variables not set in the file, unless \cmd{make -e} \\
	3. The command line       & \cmd{make CFLAGS=-O3} --- beats the Makefile      \\
	4. \texttt{override}      & \cmd{override CFLAGS += -g} --- beats even the command line \\
\end{keys}

\begin{eg}
\begin{shell}
$ make                          # CFLAGS as written in the file
$ make CFLAGS='-O0 -g'          # override for this build
$ make -e                       # let the environment win everywhere
$ make CC=clang -j8             # a different compiler, eight jobs
\end{shell}
\end{eg}

\begin{note}[Target-specific variables]
	A variable can be given a different value for one target and everything it
	depends on --- which is how a debug build is usually expressed:
\begin{makecode}
debug: CFLAGS += -O0 -g -fsanitize=address
debug: app
\end{makecode}
	\cmd{make debug} then builds \texttt{app} with the extra flags. Note the
	caveat from \autoref{sec:mk-model}: if the objects already exist from a normal
	build, make sees no changed timestamps and rebuilds nothing. Run
	\cmd{make clean} first, or give the debug build its own output directory.
\end{note}

\begin{moral}
	Every value that someone might reasonably want to change --- compiler, flags,
	install prefix, output directory --- belongs in a variable at the top of the
	file, assigned with \texttt{?=}. It costs one line and saves the next person
	from editing your recipes.
\end{moral}
